% Generated draft: requires mathematical and source review.
\documentclass[11pt]{article}
\usepackage[margin=25mm]{geometry}
\usepackage{fontspec}
\setmainfont{Latin Modern Roman}
\usepackage{amsmath,amssymb,mathtools,unicode-math}
\setmathfont{Latin Modern Math}
\usepackage{enumitem,needspace,xurl,hyperref}
\setlength{\parindent}{0pt}
\setlength{\parskip}{5pt}
\setlength{\emergencystretch}{3em}
\newcommand{\N}{\mathbb{N}}
\newcommand{\Z}{\mathbb{Z}}
\newcommand{\Q}{\mathbb{Q}}
\newcommand{\R}{\mathbb{R}}
\newcommand{\C}{\mathbb{C}}
\newcommand{\F}{\mathbb{F}}
\newcommand{\A}{\mathbb{A}}
\newcommand{\PP}{\mathbb{P}}
\newcommand{\E}{\mathbb{E}}
\newcommand{\eps}{\varepsilon}
\newcommand{\dd}{\,\mathrm{d}}
\begin{document}
\section*{UnsolvedMath ID: 11000167}
\textbf{Source problem code:} catalogue entry
\subsection{Definitions and mathematical statement}
Let $S_g$ be the genus-$g$ surface of the source, let $\mathcal M_g$ be its mapping class group, and fix a normal subgroup $G_g\triangleleft\mathcal M_g$. The proposed object $\mathcal G(S_g)$ is a complex of curves associated with this setting. The excerpt does not give a construction from $G_g$, identify its allowed vertices, or state its simplex/incidence rules.

Whenever the construction is invariant under mapping classes, the action on curves induces a natural homomorphism
\[
 \eta:\mathcal M_g\longrightarrow\operatorname{Aut}_{\mathrm{simp}}(\mathcal G(S_g)).
\]
\textbf{Problem 2.10.} Find basic structural properties of $\mathcal G(S_g)$ that guarantee that this natural map is an isomorphism. In terms of the target conclusion, every simplicial automorphism should be induced by a mapping class and the inducing mapping class should be unique. These are the desired rigidity and faithfulness conclusions, not a proposed answer by themselves: the source asks for useful properties of the complex that force them.

A precise theorem needs to specify the relationship between $G_g$ and the complex, any restrictions on the normal subgroup and genus, and the mapping-class convention. Normality alone does not define a curve complex or its incidence structure. Full curve-complex automorphism groups often detect orientation-reversing classes, so an orientation-preserving $\mathcal M_g$ and an extended mapping class group give different targets. Low-genus action kernels, in particular the usual genus-two central involution in curve actions, also obstruct a literal faithful identification unless the complex or group convention deals with them. The source supplies no such exceptions or criteria. Its final sentence returns to the article's theme rather than adding another mathematical assertion.
\subsection{Short English statement}
Identify usable conditions on a curve complex associated with a normal mapping-class subgroup that force every complex automorphism to come from a unique mapping class. The subgroup-to-complex construction, orientation convention, and low-genus exceptions are missing.
\subsection{Sources}
\begin{itemize}
\item Local UnsolvedMath registry, permanent ID 11000167, source code catalogue entry; no individual proposer or proposal year supplied.
\item Farb, edited volume \emph{Problems on Mapping Class Groups and Related Topics}, Problem 2.10, PDF page 158; \url{https://www.math.uchicago.edu/~farb/papers/mcgbook.pdf}. The registry records PDF-font-structure extraction and access on 2026-07-29; the inventory distinguishes 329 labeled headings by source page.
\item The summary relates the question to Ivanov rigidity but supplies no exact theorem or primary citation and labels its literature information truncated.
\end{itemize}
\end{document}
